\section{Keep the LaTeX source easy to reason about}
\label{sec:latex}

The source tree determines how easily authors can revise the paper. A clear
project prevents duplicated text and accidental notation changes. It also lets
several authors work without breaking a submission mode.

\subsection{Use the entry point as an outline}

Keep \texttt{main.tex} short enough to read as the manuscript's table of
contents. It should show the build mode and manuscript order at a glance. Put
prose in section files:

\begin{verbatim}
paper/
|-- main.tex
|-- references.bib
|-- body/
|   |-- abstract.tex
|   |-- introduction.tex
|   |-- method.tex
|   `-- experiments.tex
|-- appendix/
|-- figures/
`-- preamble/
\end{verbatim}

This guide preserves the complete preamble and title structure from
OGPO~\citep{patil2026ogpo} so that the original formatting remains inspectable.
A new paper should keep the modular structure but remove unused
project-specific commands. Source-specific macro sprawl makes later revisions
harder to audit. Delete duplicate package loads and commented-out manuscripts.

\subsection{Separate presentation from meaning}

Stable package and layout choices belong in a style file. Put paper-specific
notation in a commands file, including method names and operators. Keep author
comments and colored revisions in a drafting file that can be disabled before
release.

Keep venue \texttt{.sty} and \texttt{.cls} files unmodified. Load the venue
style once, before project-specific styling, and record its version. Do not
hide notation, spacing fixes, or author comments inside vendor code.

Use macros for semantic objects:

\begin{verbatim}
\newcommand{\method}{\textsc{Method}\xspace}
\newcommand{\state}{s}
\DeclareMathOperator*{\argmaxop}{arg\,max}
\end{verbatim}

A method rename should require one edit. Avoid macros that conceal a whole
sentence or arbitrary formatting. Source should remain understandable without
opening several files for each line.

\subsection{Give each toggle one responsibility}

Build toggles are useful when a release changes presentation. Use them for
anonymous and public author blocks or arXiv metadata. They can also control
acknowledgments, appendix length, and venue-approved theorem styles. They
should not select different claims or experimental values. Two build modes
with different science will drift.

Keep toggles near the top of \texttt{main.tex}. Build every supported mode
before release. The public and anonymous versions may look different, but they
should make the same scientific claims.

\subsection{Keep figure sources reproducible}

Store final vector diagrams and plots as PDF when possible, and raster task
images as PNG. Keep a documented generation path for assets produced from
code. Use stable descriptive filenames and labels. Place the figure and its
caption in the section that argues from it.

Do not tune layout with repeated negative \texttt{\textbackslash vspace}
commands while the content is moving. Settle the figure, caption, and section
order first. Local spacing patches are difficult to maintain and often break
under another venue style.

\subsection{Compile as part of writing}

Compile after each meaningful change. A practical loop is:

\begin{enumerate}[leftmargin=*]
    \item run \texttt{latexmk} with errors treated as failures;
    \item inspect undefined references, citations, missing files, and boxes
    that exceed the margin;
    \item render the PDF and inspect every page at final size;
    \item repeat after changes to figures, tables, title blocks, or build
    modes.
\end{enumerate}

Before submission, build from a clean directory. Remove TODOs, author comments,
and revision colors. Treat missing references and undefined citations as
errors. Confirm that the source bundle contains every figure and bibliography
file. A successful compiler exit is necessary, but it does not catch a clipped
label or an unreadable plot. It also cannot tell you that a page is empty or a
title crosses the margin.
